\subsection{STEP6-01--03: an exact candidate for the two-indicator Weil test}
The timeline does not define its $W$. We can nevertheless give a precise candidate normalization, identify its arithmetic cross term, and prove the stated half-plane shape. None of these calculations attributes this normalization to the historical determinant construction.

Put $g=\mathbf1_{(0,1]}$, $g_X=\mathbf1_{(0,X^{-1}]}$ and $f^\pm_{X,\eta}=g\pm X^\eta g_X$, for $X\ge1$ and $\eta\in\R$. On these real functions define
\[
 \M f(s)=\int_0^\infty f(x)x^s\,\frac{dx}{x}\quad(\Re s>0),
 \qquad W_{\rm std}(f)=\sum_\rho\M f(\rho)\M f(1-\rho),
\]
with all nontrivial zeros and their multiplicities. This is a quadratic form in $f$. If $W$ instead denotes the linear zero functional, this expression applies it to $f*_\times f^\sharp$, where $f^\sharp(x)=x^{-1}\overline{f(x^{-1})}$: its Mellin transform is $\M f(s)\overline{\M f(1-\bar s)}$. Fubini is valid for these functions on $0<\Re s<1$. For example, $g*_\times g^\sharp(x)=\min(1,x^{-1})$, not $g(x)$.

\paragraph{STEP6-01: the trace and its positive product.}
Use the paired-zero Hilbert space of \fid{B-01}, put
\[
 Ae_j=\gamma_je_j,\qquad
 D(A)=\{v:\sum_j|\gamma_j|^2|v_j|^2<\infty\},\qquad
 C=(A^2+I/4)^{-1}=B(I+B/4)^{-1}.
\]
Here $A$ is the maximal closed normal diagonal operator. For $L=\log X$ and $a=X^{\eta-1/2}$,
\begin{equation}\label{eq:WeilIndicatorTrace}
 W_{\rm std}(f^\pm_{X,\eta})
 =2\Tr\!\left[C\bigl((1+a^2)I\pm2a\cos(AL)\bigr)\right],
 \qquad W_{\rm std}(g)=2\Tr C=2\lambda_1.
\end{equation}
Under RH, $A$ is selfadjoint, $U=e^{iAL}$ is unitary and, writing $F_\pm=I\pm aU$,
\begin{equation}\label{eq:WeilIndicatorPositive}
 W_{\rm std}(f^\pm_{X,\eta})
 =2\Tr(CF_\pm^*F_\pm)
 =2\|F_\pm C^{1/2}\|_{\mathcal S_2}^{2}\ge0.
\end{equation}
\begin{proof}
Direct absolutely convergent integration gives $\M f^\pm(s)=(1\pm X^{\eta-s})/s$ for $\Re s>0$. If $\rho=\beta+it$ is a nontrivial zero, the numerator product is bounded by $(1+X^\eta)^2$ and $|\rho(1-\rho)|\ge t^2$. The zero count $N(T)=O(T\log T)$ gives
$\sum_\rho|\rho(1-\rho)|^{-1}<\infty$: the dyadic annulus $2^m\le|t|<2^{m+1}$ contributes $O((m+1)2^{-m})$, and the finitely many low zeros have nonzero denominators. Thus the quadratic sum is absolutely convergent, locally uniformly in $X,\eta$, and can be paired or conjugated.

The entries $c_j=(\gamma_j^2+1/4)^{-1}=[\rho_j(1-\rho_j)]^{-1}$ define a trace-class $C$. The sequences $c_j$ and $\gamma_j^2c_j$ are bounded: the denominators never vanish and $|\gamma_j|\to\infty$. Hence $C$ maps the Hilbert space into $D(A^2)$ and is the two-sided inverse of $A^2+I/4$ on the appropriate domains. Since $|\Im\gamma_j|<1/2$, the diagonal operators $e^{\pm iAL}$ and $\cos(AL)$ are bounded, with entry bounds $X^{1/2}$ and $\cosh(L/2)$. Their products with $C$ are trace class.

For either sign, expansion gives
\[
 \M f^\pm(\rho)\M f^\pm(1-\rho)
 =\frac{1+X^{2\eta-1}\pm X^{\eta-\rho}\pm X^{\eta-1+\rho}}
 {\rho(1-\rho)}
 =\frac{1+a^2\pm2a\cos(\gamma L)}{\gamma^2+1/4}.
\]
This is even in $\gamma$, so every reflection pair contributes twice the representative's diagonal entry. It proves the first identity in \eqref{eq:WeilIndicatorTrace}. The base test gives $2\Tr C$, and \fid{B-03} gives $\Tr C=\lambda_1$.

Under RH, $C\ge0$ and all these diagonal operators commute. Therefore
$F_\pm^*F_\pm=(1+a^2)I\pm a(U+U^*)=(1+a^2)I\pm2a\cos(AL)$.
The operator $C^{1/2}$ is Hilbert--Schmidt. Taking the trace of
$(F_\pm C^{1/2})^*(F_\pm C^{1/2})$ proves \eqref{eq:WeilIndicatorPositive}. Off RH, the exact trace remains valid, but $U^{-1}$ need not be $U^*$ and $C$ need not be positive; the positive-product step is not asserted there.
\end{proof}

\paragraph{STEP6-02: the actual smoothed sum is a weighted cosine trace.}
Let $G_\zeta(X)$ be the exact source scalar of \fid{DET-09}, namely
$\sum_{n\le X}\Lambda(n)(X/n-1)-X\log X+(1+\gamma_E)X$. Set
\[
 R(X)=\sum_\rho\frac{X^\rho}{\rho(1-\rho)},\qquad
 b(X)=\log(2\pi)-\sum_{n=1}^\infty\frac{X^{-2n}}{2n(2n+1)}.
\]
Unconditionally, for every $X\ge1$,
\begin{equation}\label{eq:GzCosine}
 G_\zeta(X)=b(X)+R(X)
 =b(X)+2\sqrt X\,\Tr\bigl(C\cos(A\log X)\bigr).
\end{equation}
All sums converge absolutely and locally uniformly. The correction satisfies
\[
 b(1)=\log(4\pi)-1,\qquad
 \log(4\pi)-1\le b(X)<\log(2\pi).
\]
Under RH, \eqref{eq:GzCosine} gives
$|G_\zeta(X)-b(X)|\le2\lambda_1\sqrt X$.
\begin{proof}
The primary explicit-formula input is the \emph{corrected} Lemma 2.1 of Ramar\'e's corrigendum \cite[pp.~1--2]{ramare}; the original article's Lemma 2.2 has sign errors. With $\psi(X)=\sum_{n\le X}\Lambda(n)$ and $\widetilde\psi(X)=\sum_{n\le X}\Lambda(n)/n$, the corrected formula is
\[
 \widetilde\psi(X)=\log X-\gamma_E+\frac{\psi(X)-X}{X}
 -\sum_\rho\frac{X^{\rho-1}}{\rho(\rho-1)}+\frac{b(X)}X.
\]
Multiply by $X$, subtract $\psi(X)$ and then subtract
$X\log X-(1+\gamma_E)X$. This proves $G_\zeta=R+b$, because the source ramp is $X\widetilde\psi-\psi$. The zero sum is absolutely convergent by $|X^\rho|\le X$ and the reciprocal-denominator estimate above. Pairing $\rho$ with $1-\rho$ gives
$X^\rho+X^{1-\rho}=2\sqrt X\cos(\gamma\log X)$ and proves the trace identity.

The corrigendum also gives
\[
 b(X)=X\int_X^\infty
 \left(\log(2\pi)+\tfrac12\log(1-t^{-2})\right)\frac{dt}{t^2}.
\]
Expand $-\log(1-t^{-2})=\sum_{n\ge1}t^{-2n}/n$ and apply Tonelli to the nonnegative series. This gives exactly the displayed series for $b$, including $X=1$, where the logarithmic singularity is integrable. The identity
$1/[2n(2n+1)]=1/(2n)-1/(2n+1)$ gives the sum $1-\log2$ at $X=1$. Positivity and monotonicity of its terms prove the bounds and continuity. For $X>1$ the equivalent elementary expression is
\[
 b(X)=\log(2\pi)+\tfrac12\log(1-X^{-2})
 +\tfrac12\left(X\log\frac{X+1}{X-1}-2\right).
\]
At prime-power endpoints the jumps of $X\widetilde\psi-\psi$ cancel, so no half-weight convention remains in $G_\zeta$. Setting $X=1$ checks
$2\lambda_1=G_\zeta(1)-b(1)=2+\gamma_E-\log(4\pi)$.
Under RH, $|\cos(\gamma_j\log X)|\le1$ and $c_j>0$, giving the asserted bound.
\end{proof}
This derivation cites the classical analytic explicit formula; it independently proves its algebraic translation, convergence, endpoint handling and trace representation. It does not construct $A$ arithmetically from a period determinant.

\paragraph{STEP6-03: exact one-sided bounds and the half-plane implication.}
The two reflected cross sums coincide, so
\begin{align}
 W_{\rm std}(f^\pm_{X,\eta})
 &=2\lambda_1(1+X^{2\eta-1})
   \pm2X^{\eta-1}\bigl(G_\zeta(X)-b(X)\bigr),\label{eq:WeilGz}\\
 W_{\rm std}(f^+_{X,\eta})\ge0
 &\Longleftrightarrow G_\zeta(X)\ge b(X)-\lambda_1(X^{1-\eta}+X^\eta),\label{eq:WeilLower}\\
 W_{\rm std}(f^-_{X,\eta})\ge0
 &\Longleftrightarrow G_\zeta(X)\le b(X)+\lambda_1(X^{1-\eta}+X^\eta).\label{eq:WeilUpper}
\end{align}
Thus the plus sign gives a lower bound, while the minus sign gives an upper bound. At $X=1$ the values are $8\lambda_1$ and $0$, consistent with $f^+=2g$ and $f^-=0$.

\begin{lemma}
If $0<\theta<1$ and either $G_\zeta(X)\ge-KX^\theta$ or $G_\zeta(X)\le KX^\theta$ for some $K\ge0$ and all sufficiently large $X$, then $\zeta$ has no zero with real part greater than $\theta$.
\end{lemma}
\begin{proof}
The scalar Mellin identity established in \fid{DET-11} is
\[
 \Phi_\zeta(s)=-\frac{\zeta'(s)/\zeta(s)}{s(s-1)}
 -\frac1{(s-1)^2}+\frac{1+\gamma_E}{s-1}\qquad(\Re s>1).
\]
It continues meromorphically to $\Re s>0$. At $s=1$, putting $h=s-1$, the expansion
$-\zeta'/\zeta=h^{-1}-\gamma_E+O(h)$ makes the first term
$h^{-2}-(1+\gamma_E)h^{-1}+O(1)$, so the polar terms cancel. At a nontrivial zero $\rho$ of multiplicity $m$, its residue is $m/[\rho(1-\rho)]\ne0$. There are no real zeros on $(0,1)$, since the positive alternating eta series equals $(1-2^{1-s})\zeta(s)$ there; for $s>1$ use the Euler product. Thus $\Phi_\zeta$ is analytic near every positive real point.

For a lower bound choose $u_0\ge0$ beyond its threshold and define $h(u)=0$ before $u_0$ and $h(u)=G_\zeta(e^u)+Ke^{\theta u}\ge0$ afterwards. The crude bound $G_\zeta(X)=O(X(1+\log X)^2)$ follows from $\Lambda(n)\le\log n$. Hence its Laplace transform $H$ converges for $\Re s>1$, where
\[
 H(s)=\Phi_\zeta(s)-\int_0^{u_0}G_\zeta(e^u)e^{-su}\,du
       +\frac{Ke^{-(s-\theta)u_0}}{s-\theta}.
\]
The compact integral is entire. This gives meromorphic continuation of $H$ to $\Re s>\theta$, analytic near every real point there, with nonzero poles at any zeta zeros there.

For clarity, the required positive-Laplace fact has a short proof. If a nonnegative locally integrable function has finite abscissa of convergence $\sigma_c$, its Laplace transform is singular at the real point $\sigma_c$. Otherwise let a disc of analytic continuation there have radius $\delta>0$. Set $\sigma_0=\sigma_c+\delta/4$ and $r=\delta/2$. The Taylor disc about $\sigma_0$ has radius at least $3\delta/4$, whereas $\sigma_0-r<\sigma_c$. Differentiation under the integral at $\sigma_0$ is justified by convergence at a slightly smaller point still above $\sigma_c$. Its Taylor series towards the left has nonnegative terms, and Tonelli gives
\[
 \sum_{n\ge0}\frac{r^n}{n!}
  \int_0^\infty u^nh(u)e^{-\sigma_0u}\,du
 =\int_0^\infty h(u)e^{-(\sigma_0-r)u}\,du<\infty,
\]
contradicting the definition of $\sigma_c$. The case $\sigma_c=-\infty$ already converges everywhere.

Here $\sigma_c\le1$. If $\sigma_c>\theta$, the displayed continuation is analytic through that real point, contradicting this fact. Thus $\sigma_c\le\theta$, so $H$ is holomorphic on $\Re s>\theta$. Meromorphic uniqueness forces $\Phi_\zeta$ to have no pole there, excluding all the indicated zeros. For an upper bound use $h(u)=Ke^{\theta u}-G_\zeta(e^u)$ beyond $u_0$; the same proof applies to $-\Phi_\zeta$ and its nonzero residues.
\end{proof}

For $0<\eta<1$, eventual nonnegativity of either specified test therefore gives zero-freeness in $\Re s>\max(\eta,1-\eta)$, by \eqref{eq:WeilLower}--\eqref{eq:WeilUpper} and boundedness of $b$. In particular $0<\eta\le1/2$ gives exactly $\Re s>1-\eta$. At $\eta=1/2$ eventual positivity of the plus-test family is RH-equivalent: the forward direction uses the lemma and reflection, and the reverse is \eqref{eq:WeilIndicatorPositive}. This proves an equivalence for the defined form, not the unprovided historical identification.

The approximate amplitude is also transparent. For any real $q(X)>0$, the same expansion gives
$W_{\rm std}(g+qg_X)=2\lambda_1(1+q^2/X)+2qR(X)/X$.
Its positivity is exactly $G_\zeta(X)\ge b(X)-\lambda_1(X/q+q)$.
If $c_1X^\eta\le q(X)\le c_2X^\eta$ eventually, for positive constants $c_1,c_2$, the same exponent follows, using $|\lambda_1|$ in the coarse constant.

The remaining source gate is concrete: specify the linear or quadratic convention for $W$, the test correspondence and normalization, and the passage from the finite period/determinant construction to this zero form. The source scalar is exactly \path{Formalism/StatementOC.lean}, lines 11--13. The actual final path in \fid{DET-09} bounds a character sum with positive weights \emph{from above on power windows}; the present lemma proves a zeta-only criterion for a global eventual bound. The source also has lower bounds on differences $G_\zeta-G_\chi$, which are a different assertion. Character averaging and window quantifiers remain additional parts of any correspondence with that path. The positive product here uses the stipulated RH; the converse lemma does not prove its own positivity premise. \fid{DET-16} below constructs the operator directly from this arithmetic scalar kernel under RH, including its full multiplicities. Its historical finite-period correspondence and a direct transfer of the original sign remain separate.
